\documentclass[aspectratio=169,10pt]{beamer}
\usetheme{Arguelles}
\usepackage{listings}
\usetikzlibrary{arrows.meta}

\title{Disjoint Set Union}
\date{\today}
\author{%
  \begin{tabular}{c}
    A \quad \small(ID 2200110) \\[4pt]
    B \quad \small(ID 2200110) \\[4pt]
    C \quad \small(ID 2200110)
  \end{tabular}%
}
\institute{Department of Computer Science and Engineering}

% Section divider slide, shown automatically whenever a new \section begins
\AtBeginSection[]{
  \begin{frame}[standout,noframenumbering]
    \Huge\insertsection
  \end{frame}
}


\newcommand{\dsuarray}[7]{%
  {\small\setlength{\tabcolsep}{3pt}
    \renewcommand{\arraystretch}{1.35}
    \begin{tabular}{r|ccccccc}
      index & 0 & 1 & 2 & 3 & 4 & 5 & 6 \\
      \hline
      \texttt{parent} & #1 & #2 & #3 & #4 & #5 & #6 & #7
    \end{tabular}
  }%
}
 
\newcommand<>{\dsucell}[1]{%
  \begingroup
    \setlength{\fboxsep}{1.5pt}%
    \setlength{\fboxrule}{0.6pt}%
    \alt#2{%
      \fcolorbox{oc-orange-6}{oc-orange-0}{\makebox[0.9em]{\strut #1}}%
    }{%
      \fcolorbox{structure.bg}{structure.bg}{\makebox[0.9em]{\strut #1}}%
    }%
  \endgroup
}

\newcommand{\dsuquadrants}[4]{%
  \begin{columns}[T,onlytextwidth]
    \begin{column}{0.46\textwidth}
      \begin{minipage}[t][2.35cm][t]{\linewidth}
        \textbf{Pseudocode}\par\smallskip
        {\small #1}
      \end{minipage}
      \par\medskip
      \uncover<3->{\textbf{Operations}}\par\smallskip
      \uncover<3->{{\small #3}}
    \end{column}
    \begin{column}{0.50\textwidth}
      \begin{minipage}[t][2.4cm][t]{\linewidth}
        \uncover<2->{%
          \textbf{Parent array}\par\medskip
          #2
        }
      \end{minipage}
      \par\medskip
      \uncover<2->{%
        \textbf{DSU Forest}\par\smallskip
        #4
      }
    \end{column}
  \end{columns}%
}

\newcommand{\dsuforest}[8][]{%
\begin{tikzpicture}[
    x=1cm, y=0.85cm,
    vertex/.style={circle, draw=oc-gray-6, fill=white, thick,
    minimum size=0.55cm, inner sep=1pt, font=\small\bfseries},
    pointer/.style={->, >=stealth, thick, draw=oc-gray-6}
  ]
  \path[use as bounding box] (-0.2,-1.3) rectangle (5.0,2.65);
  \foreach \element/\ancestor/\horizontal/\vertical in {
    0/#2/1.0/1.6, 1/#3/0.25/0.4, 2/#4/0.25/-0.8,
    3/#5/2.0/0.4, 4/#6/2.0/-0.8, 5/#7/3.4/1.6, 6/#8/4.6/1.6}{
    \ifnum\element=\ancestor\relax
    \node[vertex, draw=oc-blue-6, fill=oc-blue-0]
    (dsunode\element) at (\horizontal,\vertical) {\element};
    \else
    \node[vertex] (dsunode\element) at (\horizontal,\vertical) {\element};
    \fi
  }
  \foreach \element/\ancestor in {0/#2,1/#3,2/#4,3/#5,4/#6,5/#7,6/#8}{
    \ifnum\element=\ancestor\relax
    \draw[pointer, oc-blue-6] (dsunode\element) edge[loop above] (dsunode\element);
    \else
    \draw[pointer] (dsunode\element) -- (dsunode\ancestor);
    \fi
  }
  #1
\end{tikzpicture}%
}
 
% ---- Code-block styling for pseudocode ----
\lstdefinestyle{dsupseudocode}{
  language={},
  morekeywords={if,then,return},
  literate={<-}{{$\leftarrow$}}2 {!=}{{$\ne$}}2,
  basicstyle=\ttfamily\footnotesize\color{oc-gray-9},
  keywordstyle=\color{oc-blue-7}\bfseries,
  commentstyle=\color{oc-gray-6}\itshape,
  stringstyle=\color{oc-green-7},
  emph={make_set,find,union},
  emphstyle=\color{oc-teal-7},
  numbers=left,
  numberstyle=\tiny\color{oc-gray-6},
  numbersep=8pt,
  xleftmargin=1.5em,
  frame=none,
  columns=fullflexible,
  keepspaces=true,
  showstringspaces=false,
  breaklines=true,
  tabsize=4,
  aboveskip=4pt,
  belowskip=4pt
}

\begin{document}

\newlength{\dsucodewidth}
\setlength{\dsucodewidth}{0.46\textwidth}

\defverbatim[colored,width=\dsucodewidth]\dsumakesetcode{%
\begin{lstlisting}[style=dsupseudocode]
make_set(node):
    parent[node] <- node
\end{lstlisting}
}

\defverbatim[colored,width=\dsucodewidth]\dsufindcode{%
\begin{lstlisting}[style=dsupseudocode]
find(node):
    if parent[node] = node then
        return node
    return find(parent[node])
\end{lstlisting}
}

\defverbatim[colored,width=\dsucodewidth]\dsuunioncode{%
\begin{lstlisting}[style=dsupseudocode]
union(first, second):
    rootA <- find(first)
    rootB <- find(second)
    if rootA != rootB then
        parent[rootB] <- rootA
\end{lstlisting}
}

% ---------- Slide 1: Title ----------
\begin{frame}[plain]
  \maketitle
\end{frame}

% ---------- Slide 2: Outline ----------
\begin{frame}{Outline}
  \setbeamercolor{section in toc}{fg=oc-gray-8}
  \tableofcontents
\end{frame}

\input{intro.tex}
\input{naive_implementation.tex}
\section{Path Compression}

\tikzset{
  pc vertex/.style={circle, draw=oc-gray-6, fill=white, thick,
    minimum size=0.55cm, inner sep=1pt, font=\small\bfseries},
  pc pointer/.style={->, >=stealth, thick, draw=oc-gray-6},
  pc visited/.style={draw=oc-orange-6, fill=oc-orange-0, line width=1.5pt},
  pc shortcut/.style={pc pointer, draw=oc-teal-6, line width=1.5pt}
}

\newcommand{\pcchain}[2][]{%
  \begin{tikzpicture}[x=0.8cm, y=0.8cm]
    \path[use as bounding box] (-0.4,-0.85) rectangle (6.4,2.25);
    \foreach \element in {0,...,5}{
      \node[pc vertex] (pcnode\element) at (\element,0) {\element};
    }
    \node[pc vertex, draw=oc-blue-6, fill=oc-blue-0] (pcnode6) at (6,0) {6};
    \draw[pc pointer, oc-blue-6] (pcnode6) edge[loop below] (pcnode6);
    \foreach \element in {0,...,5}{
      \ifnum\element<#2\relax
        \pgfmathtruncatemacro{\ancestor}{\element+1}
        \draw[pc pointer] (pcnode\element) -- (pcnode\ancestor);
      \else
        \ifnum\element=5\relax
          \draw[pc pointer] (pcnode5) -- (pcnode6);
        \else
          \draw[pc shortcut] (pcnode\element) to[bend left=40] (pcnode6);
        \fi
      \fi
    }
    #1
  \end{tikzpicture}%
}

\newcommand{\pcquadrants}[4]{%
  \begin{columns}[T,onlytextwidth]
    \begin{column}{0.46\textwidth}
      \begin{minipage}[t][2.65cm][t]{\linewidth}
        \textbf{Pseudocode}\par\smallskip
        #1
      \end{minipage}
      \par\medskip
      \textbf{Operations}\par\smallskip
      {\small #3}
    \end{column}
    \begin{column}{0.50\textwidth}
      \begin{minipage}[t][2.65cm][t]{\linewidth}
        \textbf{Parent array}\par\medskip
        #2
      \end{minipage}
      \par\medskip
      \textbf{DSU Forest}\par\smallskip
      #4
    \end{column}
  \end{columns}%
}

\defverbatim[colored,width=\dsucodewidth]\pcfindcode{%
\begin{lstlisting}[style=dsupseudocode,escapeinside={(*@}{@*)}]
find(node):
    if parent[node] = node then
        return node
    (*@\textcolor{oc-teal-7}{parent[node] $\leftarrow$ find(parent[node])}@*)
    return parent[node]
\end{lstlisting}
}

\defverbatim[colored,width=\dsucodewidth]\pcdescendcode{%
\begin{lstlisting}[style=dsupseudocode,escapeinside={(*@}{@*)}]
find(node):
    if parent[node] = node then
        return node
    parent[node] <- (*@\textcolor{oc-orange-6}{find(parent[node])}@*)
    return parent[node]
\end{lstlisting}
}

\defverbatim[colored,width=\dsucodewidth]\pcrootcode{%
\begin{lstlisting}[style=dsupseudocode,escapeinside={(*@}{@*)}]
find(node):
    if parent[node] = node then
        (*@\textcolor{oc-orange-6}{return node}@*)
    parent[node] <- find(parent[node])
    return parent[node]
\end{lstlisting}
}

\defverbatim[colored,width=\dsucodewidth]\pcreturncode{%
\begin{lstlisting}[style=dsupseudocode,escapeinside={(*@}{@*)}]
find(node):
    if parent[node] = node then
        return node
    (*@\textcolor{oc-teal-7}{parent[node]}@*) <- find(parent[node])
    (*@\textcolor{oc-teal-7}{return parent[node]}@*)
\end{lstlisting}
}

\begin{frame}{Naive \texttt{find}: The Same Work Again}
  \pcquadrants
    {\dsufindcode}
    {\dsuarray{0}{0}{1}{0}{3}{5}{6}
     \par\medskip\footnotesize Exactly where the previous \texttt{union} left us.}
    {\begin{itemize}
       \item<1-> Call \texttt{find(2)}: \uncover<2->{$2\to1\to0$.}
       \item<3-> Call it again: \alert{the same two links!}
       \item<4-> The answer is returned, but no parent is updated.
     \end{itemize}
     \uncover<5->{\textcolor{oc-teal-7}{Can we remember the shortcut to $0$?}}}
    {\scalebox{0.78}{\dsuforest[{
       \draw<2->[pointer, oc-orange-6, line width=1.5pt] (dsunode2) -- (dsunode1);
       \draw<2->[pointer, oc-orange-6, line width=1.5pt] (dsunode1) -- (dsunode0);
       \draw<5->[pointer, oc-teal-6, dashed, line width=1.5pt]
         (dsunode2) to[bend left=35] (dsunode0);
     }]{0}{0}{1}{0}{3}{5}{6}}}
  \note{25 seconds, 5 overlays. Begin with the unchanged context and find(2).
    Reveal the complete two-link path, then explain that a repeat follows it again.
    Reveal the flaw, then show the question and proposed shortcut together.
    The dashed teal arrow is only a proposed shortcut; the array has not changed.}
\end{frame}

\begin{frame}{Naive Union Can Build a Chain}
  \begin{columns}[T,onlytextwidth]
    \begin{column}{0.46\textwidth}
      \textbf{A fresh worst-case example}\par\smallskip
      {\small Start again with seven singleton sets.\par
      Our union rule: \texttt{parent[rootB]} $\leftarrow$ \texttt{rootA}.}
      \medskip
      \begin{itemize}
        \item<1-> \texttt{union(1,0)}, \texttt{union(2,1)},\par
          \texttt{union(3,2)}, \ldots, \texttt{union(6,5)}.
        \item<4-> \texttt{find(0)} follows \alert{six links}.
        \item<5-> For $n$ elements: up to $n-1$ links, so one \texttt{find} can take $\Theta(n)$.
      \end{itemize}
    \end{column}
    \begin{column}{0.50\textwidth}
      \uncover<2->{\textbf{Parent array}\par\medskip
      \dsuarray{1}{2}{3}{4}{5}{6}{6}}
      \par\medskip
      \uncover<3->{\textbf{DSU Forest}\par
      \pcchain[{
        \foreach \element in {0,...,5}{
          \pgfmathtruncatemacro{\ancestor}{\element+1}
          \draw<4->[pc pointer, oc-orange-6, line width=1.5pt]
            (pcnode\element) -- (pcnode\ancestor);
        }
        \node[font=\footnotesize, text=oc-gray-6] at (3,1.0)
          {Arrows point to parents; $6$ is the root.};
      }]{6}}\par\smallskip
      \uncover<5->{{\small Naive \texttt{union} also calls \texttt{find},\par
        so it can take $\Theta(n)$ too.}}
    \end{column}
  \end{columns}
  \note{25 seconds, 5 overlays. Present the fresh example, union rule, and operation
    sequence together. Reveal the resulting array, then the completed forest.
    Explain the six-link search, then the linear bounds for find and union together.
    This shows the completed chain, not intermediate union states.
    Keep this exact array throughout the next simulation.}
\end{frame}

\begin{frame}{Path Compression: Save the Root on the Way Back}
  \begin{columns}[T,onlytextwidth]
    \begin{column}{0.46\textwidth}
      \textbf{Before: only return the root}\par\smallskip
      \dsufindcode
      \medskip
      \uncover<2->{\textbf{After: also update each visited parent}\par\smallskip
      \pcfindcode}
    \end{column}
    \begin{column}{0.50\textwidth}
      \begin{itemize}
        \item<1-> Follow parents until reaching the root, exactly as before.
        \item<2-> As recursion returns, make \textbf{every visited node} point directly to that root.
        \item<3-> Same set. Same representative. \textbf{Shorter paths.}
      \end{itemize}
      \uncover<4->{%
        \begin{tcolorbox}[colback=oc-blue-0, colframe=oc-blue-6,
          boxrule=0.8pt, arc=1.5pt, left=6pt, right=6pt, top=4pt, bottom=4pt]
          \centering\textbf{Find the root, then remember it.}\par
          \small No change to \texttt{make\_set} or the \texttt{union} pseudocode.
        \end{tcolorbox}}
    \end{column}
  \end{columns}
  \note{25 seconds, 4 overlays. Show the old listing and explain traversal.
    Reveal the new listing with its assignment explanation. Next explain unchanged
    membership, then show the takeaway and unchanged operations together.
    Store the root in every visited call, not just the starting node.}
\end{frame}

\begin{frame}{Simulating \texttt{find(0)}: First Reach the Root}
  \pcquadrants
    {\only<1-7>{\pcdescendcode}\only<8->{\pcrootcode}}
    {\dsuarray
       {\dsucell<2>{1}}{\dsucell<3>{2}}{\dsucell<4>{3}}
       {\dsucell<5>{4}}{\dsucell<6>{5}}{\dsucell<7>{6}}{\dsucell<8>{6}}
     \par\medskip\footnotesize No assignments yet: recursive calls are still pending.}
    {\begin{itemize}
       \item<1-> Begin with the chain: call \texttt{find(0)}.
       \item<2-> Calls: $0\uncover<2->{\to1}\uncover<3->{\to2}\uncover<4->{\to3}\uncover<5->{\to4}\uncover<6->{\to5}\uncover<7->{\to6}$.
       \item<8-> \texttt{parent[6] = 6}: return $6$.
     \end{itemize}
     \uncover<8->{\textcolor{oc-teal-7}{Now the suspended calls can update their parents.}}}
    {\pcchain[{
       \foreach \element in {0,...,5}{
         \pgfmathtruncatemacro{\ancestor}{\element+1}
         \pgfmathtruncatemacro{\reveal}{\element+2}
         \draw<\reveal->[pc pointer, oc-orange-6, line width=1.5pt]
           (pcnode\element) -- (pcnode\ancestor);
       }
       \node<1>[pc vertex, pc visited] at (pcnode0) {0};
       \node<8->[pc vertex, pc visited] at (pcnode6) {6};
       \node[font=\footnotesize, text=oc-orange-6] at (3,1.0)
         {\alt<8->{Root found: return $6$.}{Follow the original parent pointers.}};
     }]{6}}
  \note{30 seconds, 8 overlays. Begin at node 0. Each of the next six clicks
    follows exactly one parent link and highlights the array entry just read.
    On the final click, check the root, return 6, and announce unwinding together.
    No parent value changes on descent. The first search still pays the full cost.}
\end{frame}

\begin{frame}{Returning: Every Visited Node Gets the Shortcut}
  \pcquadrants
    {\only<1>{\pcrootcode}\only<2->{\pcreturncode}}
    {\dsuarray
       {\dsucell<7>{\alt<7->{6}{1}}}
       {\dsucell<6>{\alt<6->{6}{2}}}
       {\dsucell<5>{\alt<5->{6}{3}}}
       {\dsucell<4>{\alt<4->{6}{4}}}
       {\dsucell<3>{\alt<3->{6}{5}}}
       {\dsucell<2>{6}}{\dsucell<1>{6}}
     \par\medskip\footnotesize
     \textcolor{oc-orange-6}{Orange cell}: active parent entry.\par
     \textcolor{oc-teal-7}{Teal arrow}: a saved shortcut.}
    {\begin{minipage}[t][2.4cm][t]{\linewidth}
       \only<1>{\texttt{find(6)} returns $6$.}
       \only<2>{\texttt{parent[5]} $\leftarrow 6$; return $6$.\par
         Already points to the root: no visible change.}
       \only<3>{\texttt{parent[4]} $\leftarrow 6$; return $6$.\par
         Replace $4\to5$ with $4\to6$.}
       \only<4>{\texttt{parent[3]} $\leftarrow 6$; return $6$.\par
         Replace $3\to4$ with $3\to6$.}
       \only<5>{\texttt{parent[2]} $\leftarrow 6$; return $6$.\par
         Replace $2\to3$ with $2\to6$.}
       \only<6>{\texttt{parent[1]} $\leftarrow 6$; return $6$.\par
         Replace $1\to2$ with $1\to6$.}
       \only<7>{\texttt{parent[0]} $\leftarrow 6$; return $6$.\par
         \textbf{Done: every non-root now points to $6$.}}
       \par\medskip
       \textcolor{oc-gray-6}{Return order: $6,5,4,3,2,1,0$.}
     \end{minipage}}
    {\only<1-2>{\pcchain{6}}%
     \only<3>{\pcchain{4}}%
     \only<4>{\pcchain{3}}%
     \only<5>{\pcchain{2}}%
     \only<6>{\pcchain{1}}%
     \only<7>{\pcchain{0}}}
  \note{50 seconds, 7 overlays. Begin with the root returning 6. Advance once per
    returning node, from 5 down to 0. Explain its assignment while the array and
    matching pointer update together. Node 5 needs no rewiring. The final update
    also shows completion. Keep the nodes stationary and emphasize that every
    visited node, not just the starting node, saves the root.}
\end{frame}

\begin{frame}{The Next \texttt{find(0)}: Six Links Become One}
  \begin{columns}[T,onlytextwidth]
    \begin{column}{0.46\textwidth}
      \textbf{Before compression}\par
      \pcchain[{
        \node[font=\small\bfseries, text=oc-orange-6] at (3,1)
          {$0\to1\to2\to3\to4\to5\to6$};
      }]{6}
      \par\smallskip
      {\small \textbf{First search:} six parent links.\par
        Without compression, every repeat costs six.}
    \end{column}
    \begin{column}{0.50\textwidth}
      \uncover<2->{\textbf{After compression: the same nodes rearranged}\par
      \begin{tikzpicture}[x=0.8cm, y=0.8cm]
        \path[use as bounding box] (-0.4,-0.45) rectangle (6.4,2.65);
        \node[pc vertex, draw=oc-blue-6, fill=oc-blue-0] (pcroot) at (3,1.75) {6};
        \draw[pc pointer, oc-blue-6] (pcroot) edge[loop above] (pcroot);
        \foreach \element in {0,...,5}{
          \node[pc vertex] (pcflat\element) at (1.2*\element,0) {\element};
          \draw[pc shortcut] (pcflat\element) -- (pcroot);
        }
        \draw<3->[pc pointer, oc-orange-6, line width=1.8pt] (pcflat0) -- (pcroot);
      \end{tikzpicture}}
      \par\smallskip
      \uncover<3->{{\small \textbf{Next search:} one parent link, $0\to6$.\par
        Other visited nodes benefit too.}}
    \end{column}
  \end{columns}
  \medskip
  \uncover<4->{%
    \begin{tcolorbox}[colback=oc-blue-0, colframe=oc-blue-6,
      boxrule=0.8pt, arc=1.5pt, left=6pt, right=6pt, top=4pt, bottom=4pt]
      \centering\small For $k\geq1$ calls to \texttt{find(0)}, with \textbf{no intervening unions}:\par
      naive: $6k$ parent-link traversals\quad
      vs.\quad compression: $6+(k-1)$.
    \end{tcolorbox}}
  \note{25 seconds, 4 overlays. Present the original tree and search costs together.
    Reveal the compressed star, then highlight the next query with its one-link
    cost and the benefit to other visited nodes. Finally show the complete count
    comparison with its conditions. Count traversed links, not all reads or writes. End this part
    here; leave the explanation of union by rank entirely to member 3.}
\end{frame}

\end{document}